% Chicken Stir-Fry with Creamy Peanut Sauce — Stovetop (both books; dairy-free)
\begin{recipebox}{Chicken Stir-Fry with Creamy Peanut Sauce}{4}{30 min}

A Thai-style chicken stir-fry in a creamy, slightly sweet satay sauce made from
scratch, with a bag of frozen stir-fry veg and a little chilli heat. Medium
spice.

\begin{ingredients}
  \item[] \textbf{Chicken and vegetables}
  \item 300\,g jasmine rice, rinsed
  \item 1 tbsp oil
  \item 500\,g chicken breast or tenderloins, sliced into bite-sized strips
  \item 750\,g frozen stir-fry vegetables (e.g. Pam's International mix)
  \item 3 garlic cloves, crushed
  \item 2 tsp grated ginger
  \item[] \textbf{Peanut sauce}
  \item 100\,g smooth peanut butter
  \item 3 tbsp soy sauce
  \item 2 tbsp sweet chilli sauce
  \item 2 tbsp lime juice
  \item 1 tbsp brown sugar or honey
  \item 1 tbsp rice vinegar
  \item Sriracha or chilli flakes, to taste
  \item 150\,ml warm water
\end{ingredients}

\begin{method}
  \item \textbf{Rice.} Cook the rice in a rice cooker or as the packet says, and
        keep it warm.
  \item \textbf{Peanut sauce.} Mix the peanut butter, soy sauce, sweet chilli,
        lime juice, sugar, vinegar and chilli in a bowl. Gradually whisk in the
        warm water until smooth. It should be creamy and thick enough to coat a
        spoon.
  \item \textbf{Chicken.} Heat the oil in a large frying pan or wok over
        medium-high heat. Cook the chicken in two batches for 2--3 minutes each,
        until lightly golden (it finishes cooking later). Put it on a plate.
  \item \textbf{Vegetables.} Add the frozen vegetables straight to the hot pan (no
        need to thaw). Stir-fry over high heat for 6--9 minutes, until hot and the
        extra water has cooked off. Do it in two batches if your pan is small. Add
        the garlic and ginger for the last minute.
  \item \textbf{Bring it together.} Return the chicken, pour over the sauce and
        simmer gently for 3--4 minutes, stirring, until the chicken is cooked
        through (75\,°C in the centre) and everything is coated. Add a splash of
        water if it's too thick. Taste and add more lime, soy or chilli.
  \item \textbf{Serve} over the rice.
\end{method}

\begin{notelist}
  \item More protein: add 200\,g frozen shelled edamame with the vegetables.
  \item Richer satay: swap 75\,ml of the water for lite coconut milk.
  \item Using a 1\,kg bag of vegetables? Use it all and keep the sauce the same;
        it's just more veg.
  \item For meal prep, keep the rice separate, refrigerate promptly and eat within
        3--4 days. Reheat until piping hot.
\end{notelist}

\end{recipebox}
